\section{Related Work}
\label{sec:related_work}


Research knowledge management has been investigated from multiple
perspectives including information organization, workflow management,
semantic representation, provenance modeling, and intelligent retrieval.
However, most existing approaches address individual aspects of the
research process rather than maintaining a persistent memory of research
evolution.


\subsection{Research Knowledge Management Systems}

Traditional knowledge management approaches focus on capturing,
organizing, and sharing organizational knowledge. Davenport and Prusak
highlighted the importance of knowledge processes for creating and
transferring valuable information inside organizations
\cite{davenport1998working}. Nonaka and Takeuchi introduced the dynamic
knowledge creation model, emphasizing the transformation between tacit
and explicit knowledge
\cite{nonaka1994dynamic}.

Although these approaches provide important foundations for knowledge
management, they were not specifically designed for computational and
scientific research environments where knowledge continuously evolves
through experiments, decisions, and artifacts.


\subsection{Scientific Workflow and Research Platforms}

Scientific workflow systems have significantly improved computational
reproducibility by enabling researchers to describe and execute complex
experiments. Pegasus, for example, provides workflow planning and
execution support for large-scale scientific applications
\cite{deelman2009pegasus}. Similarly, myExperiment supports sharing and
reuse of scientific workflows and research objects
\cite{goble2008myexperiment}.

However, workflow-oriented systems mainly preserve computational
processes and execution dependencies. They provide limited support for
capturing the reasoning process behind research decisions, hypotheses,
and evolving research directions.


\subsection{Semantic Knowledge Representation and Provenance}

Semantic technologies and knowledge graphs provide mechanisms for
representing relationships between entities and enabling machine-readable
knowledge organization. The Semantic Web vision introduced by
Berners-Lee et al. established foundations for structured knowledge
representation on the Web
\cite{bernerslee2001semantic}. Recent surveys have demonstrated the
importance of knowledge graphs for integrating heterogeneous information
sources
\cite{hogan2021knowledge}.

Provenance models have also become important for improving
trustworthiness and reproducibility. The PROV model provides a formal
representation for describing relationships between entities,
activities, and agents involved in data generation
\cite{moreau2013prov}. Workflow provenance approaches further extend
these concepts to scientific processes
\cite{garijo2013workflow}.

Nevertheless, provenance approaches generally focus on tracking data
origin and transformation history rather than maintaining a complete
research memory containing ideas, decisions, and project evolution.


\subsection{Retrieval-Augmented and Generative AI Approaches}

Recent advances in retrieval-based artificial intelligence have improved
access to large-scale information repositories. Retrieval-Augmented
Generation (RAG) approaches combine external knowledge retrieval with
language models to improve generated responses
\cite{lewis2020retrieval}. Generative AI systems have further increased
the potential for intelligent assistance in research environments
\cite{park2023generative}.

Despite these advances, most current systems operate primarily on
existing documents or retrieved information. They usually lack a
persistent representation of individual research trajectories, historical
decisions, and long-term project memory.


\subsection{Research Gap}

Table~\ref{tab:related_work_comparison} summarizes the limitations of
representative approaches compared with CRME.

\begin{table}[t]
\caption{Comparison of Research Knowledge Management Approaches}
\label{tab:related_work_comparison}

\centering
\small

\begin{tabular}{lcccc}
\toprule
Approach &
Memory &
Semantic &
Provenance &
Research Evolution
\\
\midrule

Knowledge Management &
Partial &
No &
Limited &
No
\\

Workflow Systems &
Limited &
Limited &
Yes &
No
\\

Knowledge Graphs &
Partial &
Yes &
Partial &
Limited
\\

Provenance Models &
No &
Partial &
Yes &
Limited
\\

RAG-based Systems &
Document &
Yes &
No &
No
\\

CRME &
Yes &
Yes &
Yes &
Yes
\\

\bottomrule
\end{tabular}

\end{table}


Existing approaches provide valuable capabilities for storing,
executing, representing, or retrieving research information. However,
there remains a gap in systems that combine persistent research memory,
semantic organization, provenance tracking, and lifecycle-aware research
continuity in a unified framework.

CRME addresses this gap by introducing a research memory layer that
models scientific activities as connected knowledge objects and preserves
the evolution of research processes over time.


